\section{Background, from first principles}

\subsection{Sensor networks and the scheduling problem}

A wireless sensor network has many small, cheap sensors (temperature, humidity, traffic speed, air quality) and a
\emph{sink} that collects their readings. The sensors share one radio channel. If too many transmit at once, their
signals collide and are lost, so the sink coordinates who transmits when. Time is divided into \textbf{slots}; in
each slot the channel can carry at most $C$ readings. With $M$ sensors and $C<M$, the sink must choose a subset of
sensors in every slot. That choice is called \textbf{scheduling}, and the rule that makes it is the
\textbf{scheduler} or \textbf{policy}.

Example: the Intel lab has 45 sensors and we allow 6 transmissions per 5-minute slot. In any slot, 39 sensors stay
silent, and the sink has to guess what they read.

\subsection{Age of Information (AoI)}

The \textbf{Age of Information} of sensor $i$ is how many slots ago the sink last heard from it \cite{15,14}:
\[
  a_i(t+1)=\begin{cases}1 & \text{if } i \text{ was polled in slot } t,\\ a_i(t)+1 & \text{otherwise.}\end{cases}
\]
A sensor polled just now has age 1; one silent for 10 slots has age 11. AoI measures \emph{freshness}: low AoI
means recent data. The classic scheduler, \textbf{Max-AoI}, polls the $C$ sensors with the largest age. In steady
state this becomes round robin: every sensor is polled once every $M/C$ slots. Kadota \emph{et al.}\ \cite{16}
showed such age-based policies are near-optimal \emph{for AoI}.

\begin{notebox}[The weakness of AoI]
AoI looks only at \emph{when} data arrived, never at \emph{what} it says. If sensors 1 and 2 sit a metre apart in the
same room and sensor 1 just reported 24.3\,$^\circ$C, the sink already knows sensor 2 reads about 24.3\,$^\circ$C. AoI still
says sensor 2 is ``stale'' and polls it. That transmission is mostly wasted.
\end{notebox}

\subsection{Correlation}

\textbf{Spatial correlation}: nearby sensors read similar values. \textbf{Temporal correlation}: a sensor's
reading now is close to its reading 5 minutes ago. Both mean the sink can \emph{infer} silent sensors, but only if
it knows the correlation structure. We learn that structure from historical data and store it as a
\textbf{correlation graph}: nodes are sensors, and an edge joins two sensors whose unpredictable changes move
together.

\subsection{Estimating the field: the Kalman filter}

The sink keeps an estimate $\hat x_t$ of every sensor's reading plus a covariance matrix $P_t$:
\begin{itemize}
  \item $P_{ii}$ = how uncertain the sink is about sensor $i$;
  \item $P_{ij}$ = how much the uncertainties of $i$ and $j$ move together.
\end{itemize}
Each slot has two steps:
\begin{enumerate}
  \item \textbf{Predict}: move the estimate forward in time using the learned dynamics. Uncertainty grows, because
        the world changed and nobody told the sink.
  \item \textbf{Update}: when readings arrive from the polled sensors, correct the estimate. Uncertainty about
        those sensors drops sharply, and, through $P_{ij}$, uncertainty about their \emph{correlated neighbours} drops
        too.
\end{enumerate}
The \textbf{Kalman filter} \cite{20} is the mathematically optimal way to do this when the dynamics are linear and the
noise is Gaussian. Its covariance $P_t$ is the natural, correlation-aware replacement for AoI: it tells us exactly
how much the sink does \emph{not} know about each sensor, after taking neighbours into account.

\begin{keybox}
AoI says ``sensor 2 has not reported for 10 slots''. The Kalman covariance says ``but its neighbour reported 1 slot
ago, so I am almost certain what sensor 2 reads''. Our scheduler reasons with the second statement.
\end{keybox}

\subsection{Greedy uncertainty reduction}

If the sink can poll $C$ sensors, a natural rule is: pick the sensor whose reading would remove the most total
uncertainty, update $P$ as if it had been polled, pick the next one, and so on. This is \textbf{greedy variance
reduction}, the classic sensor-selection method \cite{19}, and in our project it is the \textbf{Kalman-greedy}
baseline. Its weakness is that it is \textbf{myopic}: it only optimises the current slot. It does not plan ahead
(for example, polling a sensor now because the freshness rule will otherwise force it at a worse moment later), and
it trusts its model completely, even when the model is wrong on real data.

\subsection{Graph neural networks (GNNs)}

A GNN is a neural network for data shaped like a graph. Every node (here, every sensor) starts with a feature vector
(its age, its uncertainty, and so on). In each \textbf{message-passing layer}, each node collects its neighbours'
vectors, combines them with its own through a small neural network, and produces a new vector. After a few layers,
each sensor's vector (its \textbf{embedding}) summarises its own state \emph{and} its neighbourhood \cite{21,10}.

Why a GNN suits this problem:
\begin{itemize}
  \item The data is naturally a graph (sensors and their correlations).
  \item The same weights work for any number of sensors, so one trained model runs on 45 or 325 sensors.
  \item GNNs are \textbf{permutation-equivariant}: renumbering the sensors does not change the decision.
\end{itemize}

\subsection{Reinforcement learning and PPO}

In \textbf{reinforcement learning} an \emph{agent} observes a \emph{state}, takes an \emph{action}, receives a
\emph{reward}, and the world moves to a new state. The agent learns a \emph{policy} that maximises total future
reward. Here: the state is what the sink knows, the action is which $C$ sensors to poll, and the reward is high when
the sink's estimation error is low.

There is no ``correct answer'' to copy for each slot (the best schedule depends on the future), so supervised
learning does not fit. RL learns directly from the consequence that matters: the error.

\textbf{PPO} (Proximal Policy Optimization) \cite{23} is an \textbf{actor--critic} RL algorithm:
\begin{itemize}
  \item the \textbf{actor} (our GNN) outputs a probability for each possible choice;
  \item the \textbf{critic} predicts how good a state is, which tells us whether an action did better or worse than
        expected (the \textbf{advantage});
  \item PPO increases the probability of actions with positive advantage, but \textbf{clips} each update so the
        policy never changes too much at once. This keeps training stable.
\end{itemize}
